\documentclass[12pt,a4paper]{article}

\usepackage[a4paper,text={16.5cm,25.2cm},centering]{geometry}
\usepackage{lmodern}
\usepackage{amssymb,amsmath}
\usepackage{bm}
\usepackage{graphicx}
\usepackage{microtype}
\usepackage{hyperref}
\usepackage{minted}
\setlength{\parindent}{0pt}
\setlength{\parskip}{1.2ex}

\hypersetup
       {   pdfauthor = {  },
           pdftitle={  },
           colorlinks=TRUE,
           linkcolor=black,
           citecolor=blue,
           urlcolor=blue
       }






\begin{document}



\begin{minted}[texcomments = true, mathescape, fontsize=\small, xleftmargin=0.5em]{julia}
using dn2
\end{minted}

\section{18.2.13 Gauss-Legendrove kvadrature}
Igor Nikolaj Sok, nummat 2025


\subsection{Opis problema}
Potrebno je bilo izpeljati Gaus-Legendrovo integracijsko pravilo na dveh točkah $\int_0^h f(x)\,dx = A f(x_1) + B f(x_2) + R$


Za pravilo na dveh točkah potrebujemo eksaktnost za polinom do stopnje 3


Kot lokacije vozlišč lahko vzamemo ničle 2. Legendrovega polinoma $P(x) = 0.5 ( 3x^2 - 1)$ Ničli sta torej $x1 = \frac{-1}{\sqrt{3}}$ ter $x2 = \frac{1}{\sqrt{3}}$


Uteži integriranja dobimo iz eksaktnosti integrala za konstante (na intervalu [-1, 1]) $\int_{-1}^1 1dx = 2$ $A f(x_1) + B f(x_2) = A + B = 2$ Zaradi simetričnosti; $A = B = 1$


Pravilo na [-1, 1] je torej enako $\int_{-1}^1 f(x)dx = f(\frac{-1}{\sqrt(3)}) + f(\frac{1}{\sqrt(3)})$


Napaka Gaus-Legendrovega pravila na n točkah je: $Rf = \frac{(b-a)^{2n+1} (n!)}{(2n+1)((2n)!)^3}f^{(2n)}(\xi)$


V našem primeru je napaka torej enaka $\frac{f^{(4)}(\xi)}{135}$ Ko transformiramo na poljuben interval z n podintervali je napaka potem enaka $\frac{n^5 f^{(4)}(\xi)}{4320}$


Za transformacijo pravila na interval [0, h] zamenjamo spremenljivko xx v intervalu z $\frac{h}{2}(t+1)$, kjer je t element intervala [-1, 1] dx je torej enako $\frac{h}{2}dt$


Integral torej postane: $\int_0^h f(x)dx = \frac{h}{2}(f(\frac{h}{2}(\frac{-1}{\sqrt(3)} + 1)) + f(\frac{h}{2}(\frac{1}{\sqrt(3)} + 1)))$


Če ga hočemo iz intervala [0, h] prestaviti na interval [x\emph{i, x}\{i+1\}] preprosto v funkciji prištejemo $x_i$. Velja da $x_{i+1} - x_i = h$ Integral torej postane:  $\int_{x_i}^{x_{i+1}} f(x)dx = \frac{h}{2}(f(x_i + \frac{h}{2}(\frac{-1}{\sqrt(3)} + 1)) + f(x_i + \frac{h}{2}(\frac{1}{\sqrt(3)} + 1)))$


\begin{minted}[texcomments = true, mathescape, fontsize=\small, xleftmargin=0.5em]{julia}
f(x) = sin(x) / x

interval = Interval(0, 5)
integral = Integral(f, interval)

atol = 1e-10
max_iter = 100
n = 2
prev_rez = integriraj(integral, Gaus_legendre(2))
\end{minted}
\begin{minted}[texcomments = true, mathescape, fontsize=\small, xleftmargin=0.5em, frame = leftline]{text}
1.551547590515725
\end{minted}

V primeru računanja vrednosti integrala $\int_0^5 \frac{sin(x)}{x}dx$ natančnost na 10 mest dosežemo, ko je razlika med dvema zaporednima približkoma pri različni delitni na podintervale manjša kot 1e-10 V vsakem koraku število podintervalov podvojimo ter tako dobimo oceno števila potrebnih korakov Ocenjeno število korakov za izračun integrala na 10 mest natančno je 256


\begin{minted}[texcomments = true, mathescape, fontsize=\small, xleftmargin=0.5em]{julia}
for i in 0:max_iter
    global n
    global prev_rez
    n *= 2
    rez = integriraj(integral, Gaus_legendre(n))
    if abs(prev_rez - rez) < atol
        println(rez)
        println("St korakov: ", n)
        break
    end
    println(n, " ", rez)
    println(abs(prev_rez - rez))
    prev_rez = rez
end
\end{minted}
\begin{minted}[texcomments = true, mathescape, fontsize=\small, xleftmargin=0.5em, frame = leftline]{text}
4 1.550019196137019
0.0015283943787058885
8 1.5499365643181044
8.263181891465798e-5
16 1.5499315747210947
4.98959700978574e-6
32 1.5499312655140978
3.0920699689573894e-7
64 1.5499312462296144
1.928448334354016e-8
128 1.5499312450249736
1.2046408315313784e-9
1.5499312449496934
St korakov: 256
\end{minted}


\end{document}
